{چارچوب‌محور بودن} به این معناست که معماری، ساختار پوشه‌ها، و منطق تجاری نرم‌افزار کاملا بر اساس قوانین، ساختار و امکانات یک فریمورک خاص (مانند \lr{Django، Spring، Laravel, NestJS}) شکل بگیرد. در این حالت، ابزارها و فریمورک‌ها در قلب سیستم قرار دارند و کدهای هسته برنامه به شدت به کامپوننت‌های داخلی فریمورک وابسته می‌شوند.

وابستگی شدید به یک فریمورک مشکلات جدی در طول عمر پروژه ایجاد می‌کند:
\begin{itemize}
    \item {وابستگی شدید:} منطق بیزینس مستقیما کلاس‌های فریمورک را ارث‌بری می‌کند یا از ابزارهای آن استفاده می‌کند. اگر فریمورک تغییر کند یا منسوخ شود، بخش زیادی از سیستم باید بازنویسی شود.
    \item {دشواری در تست‌نویسی:} برای نوشتن یک \lr{Unit Test} ساده روی منطق بیزینس، مجبورید ساختار فریمورک، وب‌سرور یا دیتابیس متصل به آن را شبیه‌سازی یا اجرا کنید که سرعت و دقت تست‌ها را کاهش می‌دهد.
    \item {کابوس به‌روزرسانی:} با معرفی نسخه‌های جدید فریمورک، به‌روزرسانی پروژه بسیار زمان‌بر، پرهزینه و خطرناک خواهد بود.
    \item {محدودیت در انعاف‌پذیری:} توسعه‌دهنده مجبور است کدهای خود را با ساختار پیشنهادی فریمورک هماهنگ کند، حتی اگر آن ساختار برای حل مسئله‌ی خاص بیزینس بهینه نباشد.
\end{itemize}

\lr{Clean Architecture} با تغییر جهت وابستگی‌ها و قرار دادن {منطق بیزینس در مرکز}، این مشکلات را حل می‌کند:

\begin{itemize}

\item{ قانون وابستگی}:

در این معماری، لایه‌ها به صورت دایره‌های متمرکز طراحی می‌شوند. قانون طلایی این است: \textit{وابستگی‌ها فقط باید به سمت داخل باشند.} یعنی لایه‌های درونی (بیزینس) هیچ اطلاعی از لایه‌های بیرونی (فریمورک، دیتابیس) ندارند. فریمورک صرفا یک ابزار در بیرونی‌ترین لایه است.

\item{ جداسازی منطق تجاری}:
\begin{itemize}
    \item \lr{Entities}: قوانین کلی و مستقل بیزینس را کپسوله‌سازی می‌کنند.
    \item \lr{Use Cases}: منطق اختصاصی هر عملیات سیستم را پیاده‌سازی می‌کنند.
\end{itemize}
این دو لایه مرکزی کاملا با کد خالص بدون وارد کردن پکیج‌ها یا دکوراتورهای فریمورک نوشته می‌شوند.

\item{ اصل وارونگی وابستگی}:

اگر لایه داخلی نیاز به ارتباط با یک ابزار بیرونی (مانند دیتابیس یا سرویس پیام‌رسان) داشته باشد، مستقیما به آن وصل نمی‌شود؛ بلکه یک {اینترفیس} در لایه داخلی تعریف می‌کند و فریمورک در لایه بیرونی موظف است آن را پیاده‌سازی کند.

\end{itemize}

با استفاده از \lr{Clean Architecture} {فریمورک تبدیل به یک جزئیات می‌شود}، نه محور سیستم. شما می‌توانید لایه وب یا دیتابیس فریمورک را بدون دست زدن به کدهای اصلی بیزینس تعویض یا به‌روزرسانی کنید.

